\section{Conclusion}
\label{sec:conclusion}

This paper presented a service-oriented 4D path-planning architecture for structured U-space corridor operations. The architecture separates topology provision, A*-based corridor routing, CFMS-based trajectory generation, and NDMap-based deconfliction into dedicated services. The planning pipeline first selects a corridor route, then synthesizes phase-structured waypoints, and finally generates a dense time-parameterized 4D trajectory for downstream validation and supervision.

The Hamburg vertiport corridor-network evaluation shows that A* path computation and waypoint synthesis are not the limiting stages of the evaluated planner. Both stages remain below one millisecond in the measured benchmark, while CFMS trajectory generation dominates the planner-stage runtime and largely determines the observed scaling behavior with route distance.

Future work should extend the evaluation to larger and denser corridor networks, concurrent request loads, and more operationally realistic constraints including terrain, weather, and obstacle avoidance. A further step is to evaluate the complete closed-loop interaction between trajectory generation and deconfliction, including repeated replanning triggered by detected conflicts. For the evaluated Hamburg corridor network, graph search is not the bottleneck. The main computational cost lies in CFMS-based trajectory generation, so future optimization should focus on trajectory generation, sampling, and reuse of route segments.
